\documentclass[../../main.tex]{subfiles}

\begin{document}

\section{Urysohn 引理和 Tietze 延拓定理}\label{section:ch12-Urysohn-Tietze}

在度量空间 $(X,\rho)$ 上有丰富的连续实值函数. 事实上, 对 $X$ 的非空闭子集 $C$, 称 $d_{C}$ 为距离 $C$ 的函数, 定义为
\begin{align*}
\text{对所有 }x\in X,\quad d_{C}(x)=\inf_{x'\in C}\rho(x',x)
\end{align*}
是连续的, 而 $C$ 是 $0$ 在 $d_{C}$ 下的原象. 连续性从三角不等式得到. 此外, 若 $A$ 和 $B$ 是 $X$ 的不交闭子集, 存在 $X$ 上的连续实值函数 $f$ 使得
\begin{align*}
f(X)\subseteq[0,1],\quad\text{在 }A\text{ 上 }f=0,\quad\text{在 }B\text{ 上 }f=1
\end{align*}
该函数由 $X$ 上的 $f=\dfrac{d_{A}}{d_{A}+d_{B}}$ 给出.

$f$ 的直接构造依赖于 $X$ 上的度量. 然而, 下面的基本引理告诉我们, 在任何正规的拓扑空间上, 特别地, 在任何紧的 Hausdorff 空间上, 存在这样的函数.

\begin{theorem}[Urysohn 引理]\label{theorem:ch12-Urysohn}
令 $A$ 和 $B$ 是正规拓扑空间 $X$ 的不交闭子集. 则对任何实数的闭有界区间 $[a,b]$, 存在定义在 $X$ 上的在 $[a,b]$ 中取值的连续实值函数 $f$, 满足在 $A$ 上 $f=a$ 而在 $B$ 上 $f=b$.
\end{theorem}

该引理可考虑为延拓结果: 事实上, 通过设在 $A$ 上 $f=a$ 而在 $B$ 上 $f=b$ 定义 $A\cup B$ 上的实值函数 $f$. 这是 $X$ 的闭子集 $A\cup B$ 上的连续函数, 它在 $[a,b]$ 取值. Urysohn 引理断言该函数可被延拓为整个 $X$ 上的连续函数, 它也在 $[a,b]$ 取值. 我们注意到若一个 Tychonoff 拓扑空间 $X$ 具有 Urysohn 引理描述的性质, 则 $X$ 必须是正规的. 事实上, 对 $X$ 的非空不交闭子集 $A$ 与 $B$ 以及在 $A$ 上取值 $0$ 而 $B$ 上取值 $1$ 的 $X$ 上的连续函数 $f$, 若 $I_{1}$ 和 $I_{2}$ 分别是包含 $0$ 与 $1$ 的不交开区间, 则 $f^{-1}(I_{1})$ 与 $f^{-1}(I_{2})$ 分别是 $A$ 和 $B$ 的不交邻域.

若我们引入以下概念并且证明两个预备结果, 则 Urysohn 引理的证明就更为清晰.

\begin{definition}[正规上升族]\label{definition:ch12-正规上升}
令 $X$ 为拓扑空间而 $\Lambda$ 是实数集合. 被 $\Lambda$ 标记的 $X$ 的开子集族 $\{\mathcal{O}_{\lambda}\}_{\lambda\in\Lambda}$ 称为是正规上升的, 若对任何 $\lambda_{1},\lambda_{2}\in\Lambda$,
\begin{align*}
\text{若 }\lambda_{1}<\lambda_{2},\quad\text{则 }\overline{\mathcal{O}_{\lambda_{1}}}\subseteq\mathcal{O}_{\lambda_{2}}.
\end{align*}
\end{definition}

\begin{example}\label{example:ch12-正规上升例}
令 $f$ 为拓扑空间 $X$ 上的连续实值函数. 令 $\Lambda$ 为任意实数集, 对 $\lambda\in\Lambda$, 定义
\begin{align*}
\mathcal{O}_{\lambda}=\{x\in X\mid f(x)<\lambda\}.
\end{align*}
由连续性, 显然若 $\lambda_{1}<\lambda_{2}$, 则
\begin{align*}
\mathcal{O}_{\lambda_{1}}\subseteq\{x\in X\mid f(x)\leqslant\lambda_{1}\}\subseteq\{x\in X\mid f(x)<\lambda_{2}\}=\mathcal{O}_{\lambda_{2}}.
\end{align*}
因此开集族 $\{\mathcal{O}_{\lambda}\}_{\lambda\in\Lambda}$ 是正规上升的.
\end{example}

我们把以下引理的证明留作练习.

\begin{lemma}\label{lemma:ch12-正规上升连续}
令 $X$ 为拓扑空间. 对实数的开有界区间 $(a,b)$ 的稠密子集 $\Lambda$, 令 $\{\mathcal{O}_{\lambda}\}_{\lambda\in\Lambda}$ 为 $X$ 的开子集的正规上升族. 定义 $f:X\to\mathbf{R}$ 为: 在 $X\sim\bigcup_{\lambda\in\Lambda}\mathcal{O}_{\lambda}$ 上 $f=b$, 否则设
\begin{align}
f(x)=\inf\{\lambda\in\Lambda\mid x\in\mathcal{O}_{\lambda}\}.\label{align:ch12-正规上升f}
\end{align}
则 $f:X\to[a,b]$ 是连续的.
\end{lemma}

我们接着给出前一章命题 8 的强的推广.

\begin{lemma}\label{lemma:ch12-正规延拓}
令 $X$ 为正规拓扑空间, $F$ 是 $X$ 的闭子集, 而 $\mathcal{U}$ 是 $F$ 的邻域. 则对任何开有界区间 $(a,b)$, 存在 $(a,b)$ 的稠密子集 $\Lambda$ 与 $X$ 的正规上升的开子集族 $\{\mathcal{O}_{\lambda}\}_{\lambda\in\Lambda}$, 使得对所有 $\lambda\in\Lambda$,
\begin{align}
F\subseteq\mathcal{O}_{\lambda}\subseteq\overline{\mathcal{O}_{\lambda}}\subseteq\mathcal{U}.\label{align:ch12-正规延拓}
\end{align}
\end{lemma}

\begin{proof}
由于存在把 $(0,1)$ 映上 $(a,b)$ 的严格递增连续函数, 我们可以假设 $(a,b)=(0,1)$. 对 $(0,1)$ 的稠密子集我们选取属于 $(0,1)$ 的二进有理数集:
\begin{align*}
\Lambda=\{m/2^{n}\mid m\text{ 和 }n\text{ 是自然数},1\leqslant m\leqslant 2^{n}-1\}.
\end{align*}
对每个自然数 $n$, 令 $\Lambda_{n}$ 为元素的分母为 $2^{n}$ 的 $\Lambda$ 的子集. 我们将归纳地定义正规上升的开集族 $\{\mathcal{O}_{\lambda}\}_{\lambda\in\Lambda_{n}}$, 其中每个指标是它的前一个指标的延拓.

利用前一章的命题 8, 我们可以选取开集 $\mathcal{O}_{1/2}$ 使得
\begin{align*}
F\subseteq\mathcal{O}_{1/2}\subseteq\overline{\mathcal{O}_{1/2}}\subseteq\mathcal{U}.
\end{align*}
因此我们定义了 $\{\mathcal{O}_{\lambda}\}_{\lambda\in\Lambda_{1}}$. 现在我们再用命题 8 两次, 首先是 $F$ 不变而 $\mathcal{U}=\mathcal{O}_{1/2}$, 接着 $F=\overline{\mathcal{O}_{1/2}}$ 而 $\mathcal{U}$ 不变, 找到开集 $\mathcal{O}_{1/4}$ 与 $\mathcal{O}_{3/4}$ 使得
\begin{align*}
F\subseteq\mathcal{O}_{1/4}\subseteq\overline{\mathcal{O}_{1/4}}\subseteq\mathcal{O}_{1/2}\subseteq\overline{\mathcal{O}_{1/2}}\subseteq\mathcal{O}_{3/4}\subseteq\overline{\mathcal{O}_{3/4}}\subseteq\mathcal{U}.
\end{align*}
因此我们已将正规上升族 $\{\mathcal{O}_{\lambda}\}_{\lambda\in\Lambda_{1}}$ 延拓为正规上升族 $\{\mathcal{O}_{\lambda}\}_{\lambda\in\Lambda_{2}}$. 现在对每个自然数 $n$, 如何归纳地定义开集的正规上升族 $\{\mathcal{O}_{\lambda}\}_{\lambda\in\Lambda_{n}}$ 是显然的. 观察到该可数族的并是由 $\Lambda$ 参数化的正规上升的开集族. 该族的每个集合是 $F$ 的邻域且有包含于 $\mathcal{U}$ 的紧闭包.
\end{proof}

\textbf{Urysohn 引理的证明.} 根据应用于 $F=A$ 与 $\mathcal{U}=X\sim B$ 的\cref{lemma:ch12-正规延拓}, 我们能够选取 $(a,b)$ 的稠密子集 $\Lambda$ 与 $X$ 的正规上升的开子集族 $\{\mathcal{O}_{\lambda}\}_{\lambda\in\Lambda}$, 使得
\begin{align*}
\text{对所有 }\lambda\in\Lambda,\quad A\subseteq\mathcal{O}_{\lambda}\subseteq X\sim B.
\end{align*}
定义函数 $f:X\to[a,b]$ 为在 $X\sim\bigcup_{\lambda\in\Lambda}\mathcal{O}_{\lambda}$ 上 $f=b$, 否则设
\begin{align*}
f(x)=\inf\{\lambda\in\Lambda\mid x\in\mathcal{O}_{\lambda}\}.
\end{align*}
则在 $A$ 上 $f=a$ 而在 $B$ 上 $f=b$. \cref{lemma:ch12-正规上升连续}告诉我们 $f$ 是连续的.

我们上面提到可考虑 Urysohn 引理为延拓结果. 我们现在用该引理证明强得多的延拓定理.

\begin{theorem}[Tietze 延拓定理]\label{theorem:ch12-Tietze}
令 $X$ 为正规拓扑空间, $F$ 是 $X$ 的闭子集, $f$ 是 $F$ 上取值于闭有界区间 $[a,b]$ 的连续实值函数. 则 $f$ 有到整个 $X$ 上的延拓且取值于 $[a,b]$.
\end{theorem}

\begin{proof}
由于闭有界区间 $[a,b]$ 与 $[-1,1]$ 是同胚的, 考虑 $[a,b]=[-1,1]$ 的情形是充分也是方便的. 我们通过构造 $X$ 上具有以下两个性质的连续实值函数序列 $\{g_{n}\}$ 证明: 对每个指标 $n$,
\begin{align}
\text{在 }X\text{ 上},\quad |g_{n}|\leqslant(2/3)^{n}\label{align:ch12-gn界}
\end{align}
且
\begin{align}
\text{在 }F\text{ 上},\quad |f-[g_{1}+\cdots+g_{n}]|\leqslant(2/3)^{n}.\label{align:ch12-gn逼近}
\end{align}
事实上, 暂时假定该函数序列已被构造. 对每个指标 $n$, 定义 $X$ 上的实值函数 $s_{n}$ 为
\begin{align*}
\text{对 }X\text{ 中的 }x,\quad s_{n}(x)=\sum_{k=1}^{n}g_{n}(x).
\end{align*}
我们从估计式\eqref{align:ch12-gn界}推出, 对 $X$ 中的每个 $x$, $\{s_{n}(x)\}$ 是实数的 Cauchy 序列. 由于 $\mathbf{R}$ 是完备的, 该序列收敛. 对 $X$ 中的 $x$ 定义
\begin{align*}
g(x)=\lim_{n\to\infty}s_{n}(x).
\end{align*}
由于每个 $g_{n}$ 在 $X$ 上是连续的, 每个 $s_{n}$ 也如此. 我们可以从估计式\eqref{align:ch12-gn界}推出 $\{s_{n}\}$ 在 $X$ 上一致收敛于 $g$. 因此 $g$ 是连续的. 从估计式\eqref{align:ch12-gn逼近}可知, 在 $F$ 上 $f=g$. 因此若我们构造出序列 $\{g_{n}\}$, 则定理得证. 我们用归纳法构造.

宣称: 对每个 $a>0$ 与满足在 $F$ 上 $|h|\leqslant a$ 的连续函数 $h:F\to\mathbf{R}$, 存在连续函数 $g:X\to\mathbf{R}$ 使得
\begin{align}
\text{在 }X\text{ 上 }|g|\leqslant(2/3)a\text{ 且在 }F\text{ 上 }|h-g|\leqslant(2/3)a.\label{align:ch12-逼近宣称}
\end{align}
事实上, 定义
\begin{align*}
A=\{x\text{ 属于 }F\mid h(x)\leqslant-(1/3)a\}\quad\text{与}\quad B=\{x\text{ 属于 }F\mid h(x)\geqslant(1/3)a\}.
\end{align*}
由于 $h$ 在 $F$ 上连续而 $F$ 是 $X$ 的闭子集, $A$ 和 $B$ 是 $X$ 的不交闭子集. 因此, 根据 Urysohn 引理, 存在 $X$ 上的连续实值函数 $g$ 使得
\begin{align*}
\text{在 }X\text{ 上},\quad |g|\leqslant(1/3)a,\quad g(A)=-(1/3)a\text{ 而 }g(B)=(1/3)a.
\end{align*}
显然对 $g$ 的这个选取\eqref{align:ch12-逼近宣称}成立. 运用上面的逼近, 取 $h=f$ 与 $a=1$, 可找到连续函数 $g_{1}:X\to\mathbf{R}$ 使得
\begin{align*}
\text{在 }X\text{ 上 }|g_{1}|\leqslant(2/3)\text{ 且在 }F\text{ 上 }|f-g_{1}|\leqslant(2/3).
\end{align*}
现在再次运用断言传于 $h=f-g_{1}$ 与 $a=2/3$ 以找到连续函数 $g_{2}:X\to\mathbf{R}$ 使得
\begin{align*}
\text{在 }X\text{ 上 }|g_{2}|\leqslant(2/3)^{2}\text{ 且在 }F\text{ 上 }|f-[g_{1}+g_{2}]|\leqslant(2/3)^{2}.
\end{align*}
现在如何归纳地选取具有性质\eqref{align:ch12-gn界}和\eqref{align:ch12-gn逼近}的序列 $\{g_{n}\}$ 是显然的.
\end{proof}

Tietze 延拓定理可推广到 $X$ 上的不必有界的实值函数(见习题 8).

作为 Urysohn 引理的第二个应用, 我们给出第二可数拓扑空间可度量化的充分与必要准则.

\begin{theorem}[Urysohn 度量化定理]\label{theorem:ch12-度量化}
令 $X$ 为第二可数的拓扑空间. 则 $X$ 可度量化当且仅当它是正规的.
\end{theorem}

\begin{proof}
我们已证明一个度量空间是正规的. 现在令 $X$ 为第二可数与正规的拓扑空间. 选取该拓扑的可数基 $\{\mathcal{U}_{n}\}_{n\in\mathbf{N}}$. 令 $A$ 为乘积 $\mathbf{N}\times\mathbf{N}$ 的子集. 其定义为
\begin{align*}
A=\{(n,m)\text{ 属于 }\mathbf{N}\times\mathbf{N}\mid\overline{\mathcal{U}_{n}}\subseteq\mathcal{U}_{m}\}.
\end{align*}
由于 $X$ 是正规的, 根据 Urysohn 引理, 对 $A$ 中的每对 $(m,n)$, 存在连续实值函数 $f_{n,m}:X\to[0,1]$ 使得
\begin{align*}
\text{在 }\overline{\mathcal{U}_{n}}\text{ 上 }f_{n,m}=0\text{ 而在 }X\sim\mathcal{U}_{m}\text{ 上 }f_{n,m}=1.
\end{align*}
对 $X$ 中的 $x,y$, 定义
\begin{align}
\rho(x,y)=\sum_{(n,m)\in A}\frac{1}{2^{n+m}}\left|f_{n,m}(x)-f_{n,m}(y)\right|.\label{align:ch12-度量化rho}
\end{align}
集合 $A$ 是可数的, 因此和式收敛. 不难看出这是一个度量. 我们宣称由 $\rho$ 诱导的拓扑是 $X$ 上的给定拓扑. 为证此有必要比较基. 具体地, 必须在每个点 $x\in X$ 证明以下两个性质:
\begin{enumerate}[(i)]
\item 若 $\mathcal{U}_{n}$ 包含 $x$, 则存在 $\varepsilon>0$ 使得 $B_{\rho}(x,\varepsilon)\subseteq\mathcal{U}_{n}$.
\item 对每个 $\varepsilon>0$, 存在 $\mathcal{U}_{n}$ 包含 $x$ 且 $\mathcal{U}_{n}\subseteq B_{\rho}(x,\varepsilon)$.
\end{enumerate}
我们将这些断言的证明留作练习.
\end{proof}

\begin{exercise}\label{exercise:ch12-01}
令 $C$ 为度量空间 $(X,\rho)$ 的闭子集. 证明到 $C$ 的距离函数 $d_{C}$ 是连续的且 $d_{C}(x)=0$ 当且仅当 $x$ 属于 $C$.
\end{exercise}

\begin{exercise}\label{exercise:ch12-02}
给出开区间 $(0,1)$ 上的一个连续实值函数, 它不能延拓为 $\mathbf{R}$ 上的连续函数. 这与 Tietze 延拓定理矛盾吗?
\end{exercise}

\begin{exercise}\label{exercise:ch12-03}
作为 Tietze 延拓定理的推论证出 Urysohn 引理.
\end{exercise}

\begin{exercise}\label{exercise:ch12-04}
对在 $\mathbf{R}^{n}$ 取值的函数叙述并证明 Tietze 延拓定理.
\end{exercise}

\begin{exercise}\label{exercise:ch12-05}
假定拓扑空间 $X$ 具有性质: 每个闭子集上的连续有界实值函数有到整个 $X$ 的连续延拓. 证明: 若 $X$ 是 Tychonoff 的, 则它是正规的.
\end{exercise}

\begin{exercise}\label{exercise:ch12-06}
令 $(X,\mathcal{T})$ 为正规拓扑空间而 $\mathcal{F}$ 是 $X$ 上的连续实值函数族. 证明 $\mathcal{T}$ 是由 $\mathcal{F}$ 诱导的弱拓扑.
\end{exercise}

\begin{exercise}\label{exercise:ch12-07}
证明 Urysohn 度量化定理证明中定义的函数 $\rho$ 是一个度量, 它定义的拓扑与给定的拓扑相同.
\end{exercise}

\begin{exercise}\label{exercise:ch12-08}
令 $X$ 为正规的拓扑空间, $F$ 是 $X$ 的闭子集, $f$ 是 $F$ 上的连续实值函数. 则 $f$ 有到整个 $X$ 上的连续延拓. 用以下步骤证明它:
\begin{enumerate}[(i)]
\item 应用 Tietze 延拓定理得到函数 $f\cdot(1+|f|)^{-1}:F\to[-1,1]$ 的连续延拓 $h:X\to[-1,1]$;
\item 再次应用 Tietze 延拓定理得到函数 $\phi:X\to[0,1]$, 使得在 $F$ 上 $\phi=1$ 而在 $h^{-1}(1)$ 和 $h^{-1}(-1)$ 上 $\phi=0$;
\item 考虑函数 $\overline{f}=\phi\cdot h/(1-\phi\cdot|h|)$.
\end{enumerate}
\end{exercise}

\begin{exercise}\label{exercise:ch12-09}
证明从拓扑空间 $X$ 到拓扑空间 $Y$ 的映射 $f$ 是连续的当且仅当存在 $Y$ 的拓扑子基 $\mathcal{S}$, 使得 $\mathcal{S}$ 中的每个集合在 $f$ 下的原象在 $X$ 中是开的. 用此证明: 若 $Y$ 是闭有界区间 $[a,b]$, 则 $f$ 是连续的, 当且仅当对每个实数 $c\in(a,b)$, 集合 $\{x\in X\mid f(x)<c\}$ 与 $\{x\in X\mid f(x)>c\}$ 是开的.
\end{exercise}

\begin{exercise}\label{exercise:ch12-10}
利用前一个习题证明引理 1.
\end{exercise}

\end{document}